% =====================================================================
% RESULTS --- DRAFT (regenerated from results/findings.md; combined
% TSB-StreamingAD-U + TSB-StreamingAD-M, 2026-08).
% Compile-safe fragment: no preamble; \input into main.tex in place of
% the existing \section{Results}. Requires siunitx, booktabs, graphicx,
% natbib, cleveref (already loaded by main.tex). Tables are \input from
% tables/ (paper/tables/tsb_main.tex, tsb_by_split.tex, class_d_appendix.tex).
% Every numeral traces to results/findings.md.
% =====================================================================

\section{Results}\label{sec:results}

We report the streaming evaluation of the \num{17} detectors that passed the
verification gate (\cref{sec:setup}) on both target families of the benchmark:
TSB-StreamingAD-U (\num{260} univariate datasets) and TSB-StreamingAD-M
(\num{267} multivariate datasets), \num{527} datasets in total, over \num{8668}
completed runs at seed \num{42}. This unified-protocol leaderboard is the fair
head-to-head comparison; the remaining \num{7} detectors are legitimate OOD methods
whose defining mechanism requires an adapted protocol, so they are evaluated
separately under their native arms and reported in \cref{sec:results-classd}. The
section is organised as an overall leaderboard, a per-family (U versus M) breakdown, a
per-category analysis, the global-versus-per-series normalisation dichotomy, an
efficiency analysis, a statistical-significance test, and the adapted-protocol
results. The headline is twofold and we state it plainly at the outset:
detectors that model the in-distribution feature manifold win decisively and
cheaply, and every detector that instead reads the classifier's output head or
gradients fails, several with systematically inverted scores.

\subsection{Overall ranking: a clear feature-manifold winner}\label{sec:results-overall}
\Cref{tab:tsb_main} reports the mean AUROC across all \num{527} datasets, and
\cref{fig:heatmap} visualises performance by category. The ordering is
unambiguous and the top of the table is occupied by a single, coherent family.
The eight best detectors are \emph{all} distance-, density-, or
feature-manifold estimators: Mahalanobis distance leads with mean AUROC
\num{0.826}, followed closely by DFM-PCA (\num{0.812}), and then the corrected
DiMMAD (\num{0.747}), InvAD (\num{0.740}), CatSight (\num{0.735}), M2N2
(\num{0.733}), DEEDEE (\num{0.700}), and CODiT (\num{0.684}). These eight form a
compact high-performing band spanning \numrange{0.684}{0.826}, well clear of the
\num{0.5} chance line, and there is a pronounced gap of \num{0.15} AUROC between
the weakest top-tier detector (CODiT, \num{0.684}) and the next entry. The two
leaders, Mahalanobis and DFM-PCA, are separated from the field by a comfortable
margin and are the natural default recommendation for streaming time-series OOD
detection (developed in \cref{sec:discussion}).

Below the feature-manifold tier the picture inverts sharply. A single
reconstruction detector, DiffAD (\num{0.531}), sits marginally above chance, and
then every remaining detector falls at or below the \num{0.5} line. Critically,
\emph{all} post-hoc softmax, logit, activation, and gradient detectors fail
together: MSP (\num{0.370}), DICE (\num{0.313}), SCALE (\num{0.308}), ODIN
(\num{0.303}), ReAct (\num{0.303}), energy (\num{0.301}), and GradNorm
(\num{0.294}). The seasonal-ratio detector SRS (\num{0.352}) fails alongside
them. An AUROC well below \num{0.5} is not mere weakness but \emph{inversion}:
these detectors separate ID from OOD windows with the wrong sign, assigning
higher anomaly scores to in-distribution windows than to genuinely shifted ones.
The eight post-hoc/gradient detectors cluster tightly within a \num{0.08} band
(\numrange{0.294}{0.370}), roughly half an AUROC point below the feature-manifold
leaders --- a separation between the two families so large and so consistent
that it defines the central finding of this benchmark.

% auto-generated from results/findings.md (combined TSB-StreamingAD-U + TSB-StreamingAD-M)
\begin{table}[t]\centering
\caption{Overall mean AUROC across all \num{527} TSB-StreamingAD datasets
(\num{260} univariate + \num{267} multivariate), higher is better. The top tier
(above the rule) is held exclusively by distance-, density-, and
feature-manifold detectors; every post-hoc softmax/logit/gradient detector, and
the seasonal SRS, falls below the \num{0.5} chance line. $N$ is the number of
datasets on which the detector completed.}
\label{tab:tsb_main}
\small\begin{tabular}{lrrrrr}\toprule
Method & Mean AUROC & Std & Mean FPR95 & Infer (ms) & $N$ \\\midrule
mahalanobis & 0.826 & 0.228 & 0.366 & 1.871 & 527 \\
dfm-pca & 0.812 & 0.232 & 0.390 & 3.502 & 515 \\
dimmad & 0.747 & 0.260 & 0.492 & 3.166 & 527 \\
invad & 0.740 & 0.259 & 0.490 & 2.088 & 527 \\
catsight & 0.735 & 0.295 & 0.495 & 1.675 & 515 \\
m2n2 & 0.733 & 0.276 & 0.512 & 2.171 & 527 \\
deedee & 0.700 & 0.271 & 0.512 & 1.036 & 527 \\
codit & 0.684 & 0.242 & 0.676 & 27.843 & 527 \\\midrule
diffad & 0.531 & 0.173 & 0.842 & 21.663 & 527 \\
msp & 0.370 & 0.284 & 0.877 & 1.878 & 527 \\
srs & 0.352 & 0.304 & 0.821 & 15.893 & 260 \\
dice & 0.313 & 0.305 & 0.877 & 1.756 & 527 \\
scale & 0.308 & 0.303 & 0.878 & 1.843 & 527 \\
odin & 0.303 & 0.299 & 0.884 & 8.016 & 527 \\
react & 0.303 & 0.299 & 0.886 & 1.810 & 527 \\
energy & 0.301 & 0.301 & 0.886 & 1.949 & 527 \\
gradnorm & 0.294 & 0.301 & 0.886 & 11.790 & 527 \\
\bottomrule\end{tabular}\end{table}


\begin{figure}[t]
\centering
\includegraphics[width=0.95\textwidth]{tsb_heatmap_category.png}
\caption{Mean AUROC of every detector (rows, ordered by overall combined mean)
on the three TSB-StreamingAD categories DRIFT, OOD, and STABLE (columns).
Brighter cells denote higher AUROC. The top rows --- feature distance/density and
feature-manifold detectors --- are bright across all three categories, while the
bottom rows --- post-hoc softmax, logit, activation, and gradient detectors,
together with the seasonal SRS --- are uniformly dark, indicating scores at or
below the \num{0.5} chance line. Among the leaders, the STABLE specificity
controls and the subtle DRIFT regime are separated best (Mahalanobis \num{0.855}
STABLE, \num{0.858} DRIFT; DFM-PCA \num{0.849} STABLE), while the post-hoc cluster
is darkest on STABLE, where it over-flags the very control windows a calibrated
detector should leave alone.}
\label{fig:heatmap}
\end{figure}

\subsection{Univariate and multivariate families}\label{sec:results-split}
\Cref{tab:tsb_split} reports the ranking on each family separately, and
\cref{fig:splitbar} contrasts the winning and failing clusters directly. The
leading feature-manifold cluster is stable across both regimes: Mahalanobis and
DFM-PCA are first and second on both TSB-U (\num{0.865}, \num{0.852}) and TSB-M
(\num{0.788}, \num{0.774}), and all eight top-tier detectors remain above
\num{0.59} on both families. The univariate family is somewhat easier for the
feature leaders --- Mahalanobis, DFM-PCA, DiMMAD, and CatSight each score higher
on TSB-U than on TSB-M --- with the largest family gap belonging to DEEDEE, which
drops from \num{0.810} on TSB-U to \num{0.592} on TSB-M, consistent with its
adjacent-dimension heuristic (\cref{sec:discussion}) transferring poorly to genuinely
multivariate channels. InvAD is the one leader that is marginally stronger on the
multivariate family (\num{0.700} versus \num{0.781} on U but rising to third place
on M). The post-hoc cluster remains inverted on both families, though it is
slightly \emph{less} inverted on the multivariate family (e.g.\ MSP \num{0.433} on
M versus \num{0.305} on U; SCALE \num{0.359} versus \num{0.256}), a partial
recovery that never reaches chance. There is no overlap between the two clusters
in either family: no post-hoc detector approaches even the weakest feature-manifold
detector.

% auto-generated from results/findings.md (per-split rankings)
\begin{table}[t]\centering
\caption{Mean AUROC by benchmark family: TSB-StreamingAD-U (\num{260}
univariate datasets) versus TSB-StreamingAD-M (\num{267} multivariate datasets).
Rows are ordered by the overall combined ranking of \cref{tab:tsb_main}. The
feature-manifold leaders dominate both families; the post-hoc cluster is inverted
in both, though marginally less so on the multivariate family. SRS is
univariate-seasonal by construction and its multivariate run is pending
(\emph{pend.}), so it is omitted from the multivariate column.}
\label{tab:tsb_split}
\small\begin{tabular}{lrr}\toprule
Method & AUROC (TSB-U) & AUROC (TSB-M) \\\midrule
mahalanobis & 0.865 & 0.788 \\
dfm-pca & 0.852 & 0.774 \\
dimmad & 0.798 & 0.697 \\
invad & 0.781 & 0.700 \\
catsight & 0.787 & 0.686 \\
m2n2 & 0.773 & 0.695 \\
deedee & 0.810 & 0.592 \\
codit & 0.751 & 0.618 \\\midrule
diffad & 0.523 & 0.539 \\
msp & 0.305 & 0.433 \\
srs & 0.352 & \emph{pend.} \\
dice & 0.291 & 0.333 \\
scale & 0.256 & 0.359 \\
odin & 0.256 & 0.349 \\
react & 0.253 & 0.351 \\
energy & 0.253 & 0.348 \\
gradnorm & 0.241 & 0.345 \\
\bottomrule\end{tabular}\end{table}


\begin{figure}[t]
\centering
\includegraphics[width=0.85\textwidth]{tsb_split_bar.png}
\caption{Mean AUROC per detector on TSB-StreamingAD-U versus TSB-StreamingAD-M.
The feature distance/density and feature-manifold detectors (upper band) cluster
tightly above \num{0.59} on both families, whereas the post-hoc softmax, logit,
activation, and gradient detectors (lower band) fall below the \num{0.5} chance
line on both, several with inverted scores. The two clusters are cleanly separated
with no overlap; the post-hoc cluster is marginally less inverted on the
multivariate family.}
\label{fig:splitbar}
\end{figure}

\subsection{Per-category behaviour and actionable guidance}\label{sec:results-category}
The category breakdown in \cref{fig:heatmap} turns the difficulty axis designed
into the benchmark into practical guidance. For the feature-manifold leaders, the
near-stationary STABLE controls and the subtle near-distribution DRIFT regime are
handled best --- Mahalanobis reaches \num{0.855} on STABLE and \num{0.858} on
DRIFT, DFM-PCA \num{0.849} and \num{0.825}, and DiMMAD peaks at \num{0.801} on
STABLE --- demonstrating both strong false-positive control on the specificity
controls and clean separation of a merely drifted second source. The
strong-shift OOD category is separated well but is, perhaps counterintuitively,
the most stressful for the leaders (Mahalanobis \num{0.761}, DFM-PCA \num{0.762},
InvAD \num{0.709}), indicating that a clearly different S2 regime can push feature
representations into a region where the tied-covariance model is slightly less
calibrated than on drifted or stable data. The practical reading is favourable:
the feature-manifold family is dependable across all three regimes and especially
trustworthy on the STABLE controls, so it will not raise spurious alarms on
near-stationary streams. The inverted post-hoc detectors show the mirror-image
pattern: they are darkest on STABLE (energy \num{0.260}, GradNorm \num{0.255},
ODIN \num{0.266}, ReAct \num{0.261}, SCALE \num{0.276}), meaning they most
aggressively over-flag exactly the control windows a well-calibrated detector
should ignore; MSP is least-worst on the strong-shift OOD category (\num{0.403})
but still inverted.

\subsection{The normalisation dichotomy}\label{sec:results-norm}
\Cref{fig:norm} reports mean AUROC under global versus per-series normalisation,
and the split is dramatic and directional. The feature-manifold methods peak under
\emph{global} normalisation, which preserves the second source's amplitude and
level gap: Mahalanobis reaches \num{0.867} global versus \num{0.755} per-series,
DFM-PCA \num{0.863} versus \num{0.726}, DiMMAD \num{0.832} versus \num{0.600},
InvAD \num{0.827} versus \num{0.588}, and CatSight \num{0.797} versus \num{0.630}.
The post-hoc methods invert under \emph{both} schemes but are, strikingly,
\emph{less} inverted under per-series normalisation --- the mirror image of the
feature methods: MSP \num{0.309} global against \num{0.475} per-series, energy
\num{0.214} against \num{0.452}, ODIN \num{0.218} against \num{0.451}, and
GradNorm \num{0.213} against \num{0.433}. The two clusters move in opposite
directions under the same preprocessing knob, a genuine dichotomy rather than a
shared preference. This is the empirical signature of the overconfidence mechanism
analysed in \cref{sec:discussion}: global normalisation preserves the level gap
that lets the feature geometry separate cleanly while driving the output head
further into overconfident inversion. The actionable corollary is that the winning
feature methods should be paired with global (Source-1) normalisation for
amplitude/level shifts, which is where they are strongest.

\begin{figure}[t]
\centering
\includegraphics[width=0.85\textwidth]{tsb_norm_dichotomy.png}
\caption{Mean AUROC under global (Source-1) normalisation versus per-series
normalisation, per detector. Feature distance/density detectors (upper points)
gain sharply under global normalisation; post-hoc softmax/logit/gradient detectors
(lower points) remain inverted under both schemes and are slightly less inverted
under per-series normalisation. The two clusters move in opposite directions --- a
dichotomy, not a shared preprocessing preference.}
\label{fig:norm}
\end{figure}

\subsection{Efficiency: accurate and cheap}\label{sec:results-efficiency}
\Cref{fig:efficiency} plots mean AUROC against per-sample inference time (per
\cref{tab:tsb_main}), and the winning family is Pareto-optimal: it is
simultaneously the most accurate and among the cheapest. Mahalanobis attains the
top AUROC (\num{0.826}) at just \SI{1.871}{\milli\second} per window, DFM-PCA
reaches \num{0.812} at \SI{3.502}{\milli\second}, and DEEDEE is the fastest
detector in the benchmark (\SI{1.036}{\milli\second}) while still ranking seventh
(\num{0.700}); CatSight (\SI{1.675}{\milli\second}) and InvAD
(\SI{2.088}{\milli\second}) are likewise both accurate and inexpensive. By
contrast, the three slowest detectors are all uncompetitive on accuracy: CODiT
(\SI{27.843}{\milli\second}, \num{0.684}), DiffAD (\SI{21.663}{\milli\second},
\num{0.531}), and SRS (\SI{15.893}{\milli\second}, \num{0.352}) pay an
order-of-magnitude latency premium without a corresponding accuracy gain, and
GradNorm (\SI{11.790}{\milli\second}) and ODIN (\SI{8.016}{\milli\second}) are
both slow \emph{and} inverted. For streaming deployment, where accuracy and latency
matter jointly, the feature-manifold detectors are therefore the clear default:
they hold the accuracy--efficiency front at roughly \SIrange{1}{4}{\milli\second}
per window, so the usual accuracy--latency tension does not arise. They also carry
no auxiliary-data or retraining cost, being fit on ID features alone.

\begin{figure}[t]
\centering
\includegraphics[width=0.8\textwidth]{tsb_efficiency.png}
\caption{Accuracy versus inference cost (time axis on log scale). Each point is a
detector; the upper-left region is best. Mahalanobis, DFM-PCA, DEEDEE, and CatSight
hold the Pareto front, attaining the highest mean AUROC at roughly
\SIrange{1}{4}{\milli\second} per window, whereas CODiT
($\sim$\SI{28}{\milli\second}), DiffAD ($\sim$\SI{22}{\milli\second}), and SRS
($\sim$\SI{16}{\milli\second}) are far slower without a corresponding accuracy
gain.}
\label{fig:efficiency}
\end{figure}

\subsection{Statistical significance}\label{sec:results-stats}
A Friedman test over the detectors is very highly significant:
$\chi^2 = \num{1674.97}$, $p < \num{1e-16}$ ($k = \num{17}$ detectors,
$N = \num{250}$ datasets), so the ranking differences are not attributable to
chance at any conventional level \citep{Demsar2006}. The magnitude of the
statistic, orders of magnitude beyond the $\alpha = 0.05$ threshold, reflects the
extreme and consistent spread between the feature-manifold leaders and the
inverted post-hoc cluster visible in \cref{tab:tsb_main} and \cref{fig:heatmap}.
We state one caveat honestly. This complete-case test admits only datasets on which
\emph{all} \num{17} detectors produced a score; because SRS is univariate-seasonal
and is not run on the multivariate split, that pool is $N = \num{250}$ univariate
datasets. To obtain a combined-family test over the full corpus we repeat the
Friedman test on the \num{16} detectors excluding SRS, which retains every
multivariate dataset: it is likewise very highly significant, $\chi^2 =
\num{2336.59}$, $p < \num{1e-16}$ ($k = \num{16}$, $N = \num{515}$ datasets,
univariate and multivariate). The two tests agree, so the ranking is significant
both univariate-complete and across the combined corpus.

\subsection{Adapted-protocol results}\label{sec:results-classd}
Seven detectors are legitimate OOD methods whose defining mechanism requires an
adapted protocol rather than the unified streaming protocol --- per-window shuffled
evaluation, a frozen ResNet backbone, and no auxiliary outlier corpus --- so they are
evaluated separately under their native, method-appropriate arms rather than forced
onto an implementation that does not match the source paper (\cref{tab:classd}). Three
augmented-training methods (Outlier Exposure, DivOE, DiverseMix) train against an
auxiliary outlier set that the unified protocol does not provide; two temporal/drift
methods (AE-ADWIN-LSTM, DIVERSIFY) require an ordered stream or a retrained feature
extractor; and two (DriftLens, TD-IVDM) are natively batch/window-level and, under a
per-sample reduction, coincide with feature-space detectors already present in the
unified set. The reason and native arm for each is recorded in \cref{tab:classd}; the
audit narrative behind these verdicts is developed in \cref{sec:discussion} and
\cref{app:fidelity}.

We ran these seven on both TSB-StreamingAD-U and -M under the evaluation arm most
appropriate to each and report the outcome in \cref{tab:class_d_appendix}. The result
reinforces rather than complicates the headline: no detector in this group rises
clearly above chance and none is a winner. Their per-arm AUROC ranges from \num{0.449}
to \num{0.696}, with the strongest arm (TD-IVDM as generic KDE density, \num{0.696})
still well below the feature-manifold leaders and the weakest (DivOE, head-only arm,
\num{0.449}) below chance; the augmented-training methods hover near
\num{0.45}--\num{0.59} depending on whether the head only or the full network is
fine-tuned. Evaluated even under their own favourable protocols, these methods do not
challenge the feature-manifold family, and the two that reduce to feature-space
distance/density (DriftLens, TD-IVDM) merely re-confirm the winning principle under
other names.

% auto-generated from results/findings.md (combined TSB-StreamingAD-U + TSB-StreamingAD-M class-D arms)
\begin{table}[t]\centering
\caption{Adapted-protocol detectors, evaluated under their own
method-appropriate arms on the combined TSB-StreamingAD-U and -M corpus. These
are legitimate OOD methods whose defining mechanism requires an adapted protocol
rather than the unified frozen-backbone protocol (\cref{tab:classd}), so they are
evaluated separately for fairness; reported here under the arm most favourable to
each, no detector rises clearly above chance and none is a winner.}
\label{tab:class_d_appendix}
\small\begin{tabular}{llrrr}\toprule
Method & Arm & Mean AUROC & Std & $n$ \\\midrule
tdivdm & none & 0.696 & 0.180 & 330 \\
diversemix & head-only & 0.615 & 0.253 & 327 \\
ae-adwin-lstm & none & 0.572 & 0.198 & 324 \\
diversify & cosine-centroid & 0.566 & 0.277 & 327 \\
diversemix & full-net & 0.562 & 0.249 & 327 \\
outlier-exposure & full-net & 0.560 & 0.236 & 327 \\
driftlens & none & 0.543 & 0.337 & 313 \\
divoe & full-net & 0.538 & 0.263 & 327 \\
diversify & energy & 0.493 & 0.298 & 327 \\
outlier-exposure & head-only & 0.460 & 0.262 & 327 \\
divoe & head-only & 0.449 & 0.267 & 327 \\
\bottomrule\end{tabular}\end{table}


\begin{table}[t]
\centering
\caption{The seven OOD methods evaluated under adapted, method-appropriate protocols
rather than the unified streaming protocol. Each is audited per-method
(\cref{app:fidelity}); the last column states why an adapted protocol is required and
the native arm under which the method is evaluated.}
\label{tab:classd}
\small
\begin{tabular}{l p{6.6cm} p{4.1cm}}
\toprule
Detector & Why an adapted protocol is required & Native evaluation arm \\
\midrule
Outlier Exposure & OE trains on an auxiliary outlier set; none in unified protocol & Auxiliary-training arm (else $=$ Energy) \\
DivOE & Needs auxiliary outliers + PGD synthesis + training & Auxiliary-training arm (else $=$ Energy) \\
DiverseMix & Needs auxiliary outliers + end-to-end training; at chance either orientation & Auxiliary-training arm; negative result \\
DriftLens & Native score is batch/window-level Fr\'echet drift; per-sample $=$ PCA-Mahalanobis & Batch-level drift arm \\
TD-IVDM & Multi-scale drift pillars need ordered windows; per-sample $=$ generic KDE density & Ordered multi-scale arm \\
AE-ADWIN-LSTM & LSTM + ADWIN need an ordered stream; unified eval shuffles windows & Ordered-stream arm \\
DIVERSIFY & Adversarial representation learning retrains the extractor & Learned-representation arm \\
\bottomrule
\end{tabular}
\end{table}
